\documentclass[a4paper,11pt]{article}

\usepackage[a4paper,margin=2cm]{geometry}
\usepackage[T1]{fontenc}
\usepackage[utf8]{inputenc}
\usepackage[ngerman,english]{babel}
\usepackage{lmodern}
\usepackage{microtype}
\usepackage{xcolor}
\usepackage{parskip}
\usepackage{enumitem}
\usepackage[most]{tcolorbox}
\usepackage{titlesec}
\usepackage[hidelinks]{hyperref}

\definecolor{HeaderBlueDark}{RGB}{70,110,170}
\definecolor{RowBlueLight}{RGB}{235,243,255}
\definecolor{RowBlueDark}{RGB}{220,233,250}
\definecolor{SoftGray}{RGB}{90,90,90}

\setlength{\parindent}{0pt}
\setlist[itemize]{leftmargin=1.4em,itemsep=0.35em,topsep=0.35em}
\linespread{1.06}

\titleformat{\section}[block]
{\Large\bfseries\color{HeaderBlueDark}}
{}{0pt}{}
\titlespacing{\section}{0pt}{1.3em}{0.8em}

\titleformat{\subsection}[block]
{\large\bfseries\color{HeaderBlueDark}}
{}{0pt}{}
\titlespacing{\subsection}{0pt}{1.0em}{0.6em}

\newtcolorbox{exbox}{enhanced,breakable,colback=RowBlueLight,colframe=RowBlueDark,boxrule=0.4pt,arc=1.5mm,left=1.2mm,right=1.2mm,top=1mm,bottom=1mm,title={Beispiele \textnormal{(Examples)}},fonttitle=\bfseries,colbacktitle=RowBlueDark,coltitle=black}

\NewDocumentEnvironment{grammarentry}{mm+b}{%
    \begin{tcolorbox}[enhanced,breakable,colback=white,colframe=HeaderBlueDark,boxrule=0.6pt,arc=1.5mm,left=1.5mm,right=1.5mm,top=1mm,bottom=1mm,colbacktitle=HeaderBlueDark,coltitle=white,fonttitle=\bfseries,title={#1\hfill\normalfont\footnotesize #2}]
    #3
    \end{tcolorbox}
}{}

\newcommand{\GermanText}[1]{\textbf{Deutsch: }#1\par}
\newcommand{\EnglishText}[1]{{\small\color{SoftGray}\textbf{English: }#1\par}}
\newcommand{\ExampleItem}[2]{\item \textbf{#1}\\{\small\color{SoftGray}#2}}

\title{Die Konjunktion \glqq sondern\grqq}
\date{}

\begin{document}
    \maketitle
    \tableofcontents
    \newpage

\section{Grundbedeutung}

\begin{grammarentry}{sondern}{coordinating conjunction / nebenordnende Konjunktion}
\GermanText{\textbf{sondern} benutzt man, wenn der erste Teil verneint ist und der zweite Teil die richtige Alternative zeigt.}
\EnglishText{\textbf{sondern} is used when the first part is negated and the second part gives the correct alternative.}

\begin{exbox}
\begin{itemize}
    \ExampleItem{Ich trinke nicht Kaffee, sondern Tee.}{I do not drink coffee, but tea.}
    \ExampleItem{Er wohnt nicht in Wien, sondern in Graz.}{He does not live in Vienna, but in Graz.}
    \ExampleItem{Das ist kein Fehler, sondern eine Regel.}{That is not a mistake, but a rule.}
\end{itemize}
\end{exbox}
\end{grammarentry}

\section{Die wichtigste Struktur}

\begin{grammarentry}{nicht A, sondern B}{not A, but B}
\GermanText{Die normale Struktur ist: \textbf{nicht A, sondern B}. Der zweite Teil ersetzt den ersten Teil.}
\EnglishText{The normal structure is: \textbf{not A, but B}. The second part replaces the first part.}

\begin{exbox}
\begin{itemize}
    \ExampleItem{Ich komme nicht heute, sondern morgen.}{I am not coming today, but tomorrow.}
    \ExampleItem{Sie kauft nicht ein Auto, sondern ein Fahrrad.}{She is not buying a car, but a bicycle.}
    \ExampleItem{Wir sprechen nicht über Politik, sondern über Grammatik.}{We are not talking about politics, but about grammar.}
\end{itemize}
\end{exbox}
\end{grammarentry}

\begin{grammarentry}{kein A, sondern B}{no / not a A, but B}
\GermanText{Auch \textbf{kein} ist eine Verneinung. Deshalb kann danach \textbf{sondern} stehen.}
\EnglishText{\textbf{kein} is also a negation. Therefore, \textbf{sondern} can follow it.}

\begin{exbox}
\begin{itemize}
    \ExampleItem{Ich habe keinen Bruder, sondern eine Schwester.}{I do not have a brother, but a sister.}
    \ExampleItem{Das ist keine Frage, sondern eine Antwort.}{That is not a question, but an answer.}
    \ExampleItem{Er braucht kein Taxi, sondern einen Bus.}{He does not need a taxi, but a bus.}
\end{itemize}
\end{exbox}
\end{grammarentry}

\section{Wortstellung}

\begin{grammarentry}{sondern und Wortstellung}{word order}
\GermanText{\textbf{sondern} ist eine nebenordnende Konjunktion. Es verbindet zwei gleiche Teile. Es verändert die Wortstellung nicht.}
\EnglishText{\textbf{sondern} is a coordinating conjunction. It connects two equal parts. It does not change the word order.}

\begin{exbox}
\begin{itemize}
    \ExampleItem{Ich gehe nicht nach Hause, sondern ich bleibe hier.}{I am not going home, but I am staying here.}
    \ExampleItem{Er lernt nicht Deutsch, sondern er lernt Englisch.}{He is not learning German, but he is learning English.}
    \ExampleItem{Wir fahren nicht mit dem Zug, sondern wir fahren mit dem Bus.}{We are not going by train, but we are going by bus.}
\end{itemize}
\end{exbox}
\end{grammarentry}

\begin{grammarentry}{Ellipse mit sondern}{omitting repeated words}
\GermanText{Wenn Wörter im zweiten Teil gleich sind, kann man sie oft weglassen.}
\EnglishText{If words in the second part are the same, they can often be omitted.}

\begin{exbox}
\begin{itemize}
    \ExampleItem{Ich gehe nicht nach Hause, sondern ins Büro.}{I am not going home, but to the office.}
    \ExampleItem{Er lernt nicht Deutsch, sondern Englisch.}{He is not learning German, but English.}
    \ExampleItem{Wir fahren nicht mit dem Zug, sondern mit dem Bus.}{We are not going by train, but by bus.}
\end{itemize}
\end{exbox}
\end{grammarentry}

\section{Komma}

\begin{grammarentry}{Komma vor sondern}{comma before sondern}
\GermanText{Vor \textbf{sondern} steht normalerweise ein Komma.}
\EnglishText{There is normally a comma before \textbf{sondern}.}

\begin{exbox}
\begin{itemize}
    \ExampleItem{Ich will nicht schlafen, sondern arbeiten.}{I do not want to sleep, but to work.}
    \ExampleItem{Das Buch ist nicht teuer, sondern billig.}{The book is not expensive, but cheap.}
    \ExampleItem{Sie ist nicht müde, sondern krank.}{She is not tired, but sick.}
\end{itemize}
\end{exbox}
\end{grammarentry}

\section{sondern oder aber?}

\begin{grammarentry}{sondern}{correction after negation}
\GermanText{Man benutzt \textbf{sondern}, wenn der erste Teil verneint ist und der zweite Teil die richtige Alternative ist.}
\EnglishText{Use \textbf{sondern} when the first part is negated and the second part is the correct alternative.}

\begin{exbox}
\begin{itemize}
    \ExampleItem{Ich komme nicht am Montag, sondern am Dienstag.}{I am not coming on Monday, but on Tuesday.}
    \ExampleItem{Das ist nicht mein Handy, sondern dein Handy.}{That is not my phone, but your phone.}
\end{itemize}
\end{exbox}
\end{grammarentry}

\begin{grammarentry}{aber}{contrast}
\GermanText{Man benutzt \textbf{aber}, wenn es nur einen Gegensatz gibt. Der erste Teil muss nicht falsch sein.}
\EnglishText{Use \textbf{aber} when there is only a contrast. The first part does not have to be wrong.}

\begin{exbox}
\begin{itemize}
    \ExampleItem{Ich bin müde, aber ich arbeite weiter.}{I am tired, but I keep working.}
    \ExampleItem{Das Auto ist teuer, aber gut.}{The car is expensive, but good.}
\end{itemize}
\end{exbox}
\end{grammarentry}

\begin{grammarentry}{typischer Unterschied}{typical difference}
\GermanText{\textbf{aber} zeigt oft einen Gegensatz. \textbf{sondern} korrigiert oft eine falsche Alternative nach einer Verneinung.}
\EnglishText{\textbf{aber} often shows a contrast. \textbf{sondern} often corrects a wrong alternative after a negation.}

\begin{exbox}
\begin{itemize}
    \ExampleItem{Ich bin müde, aber glücklich.}{I am tired, but happy.}
    \ExampleItem{Ich bin nicht müde, sondern krank.}{I am not tired, but sick.}
\end{itemize}
\end{exbox}
\end{grammarentry}

\section{nicht nur \ldots{}, sondern auch}

\begin{grammarentry}{nicht nur A, sondern auch B}{not only A, but also B}
\GermanText{Die Struktur \textbf{nicht nur A, sondern auch B} bedeutet: A ist richtig und B ist auch richtig.}
\EnglishText{The structure \textbf{not only A, but also B} means: A is true and B is also true.}

\begin{exbox}
\begin{itemize}
    \ExampleItem{Er spricht nicht nur Deutsch, sondern auch Englisch.}{He speaks not only German, but also English.}
    \ExampleItem{Sie arbeitet nicht nur schnell, sondern auch genau.}{She works not only fast, but also accurately.}
    \ExampleItem{Das ist nicht nur wichtig, sondern auch dringend.}{That is not only important, but also urgent.}
\end{itemize}
\end{exbox}
\end{grammarentry}

\section{Typische Fehler}

\begin{grammarentry}{Fehler 1: sondern ohne Verneinung}{wrong use without negation}
\GermanText{\textbf{sondern} braucht normalerweise eine Verneinung im ersten Teil. Ohne Verneinung benutzt man meistens \textbf{aber}.}
\EnglishText{\textbf{sondern} normally needs a negation in the first part. Without negation, one usually uses \textbf{aber}.}

\begin{exbox}
\begin{itemize}
    \ExampleItem{Falsch: Ich bin müde, sondern ich arbeite.}{Wrong: I am tired, but I work.}
    \ExampleItem{Richtig: Ich bin müde, aber ich arbeite.}{Correct: I am tired, but I work.}
\end{itemize}
\end{exbox}
\end{grammarentry}

\begin{grammarentry}{Fehler 2: falsche Wortstellung}{wrong word order}
\GermanText{Nach \textbf{sondern} kommt keine Inversion wie nach manchen Adverbien.}
\EnglishText{After \textbf{sondern}, there is no inversion as after some adverbs.}

\begin{exbox}
\begin{itemize}
    \ExampleItem{Falsch: Ich gehe nicht nach Hause, sondern bleibe ich hier.}{Wrong: I am not going home, but stay I here.}
    \ExampleItem{Richtig: Ich gehe nicht nach Hause, sondern ich bleibe hier.}{Correct: I am not going home, but I am staying here.}
    \ExampleItem{Auch richtig: Ich gehe nicht nach Hause, sondern bleibe hier.}{Also correct: I am not going home, but stay here.}
\end{itemize}
\end{exbox}
\end{grammarentry}

\section{Kurze Zusammenfassung}

\begin{grammarentry}{sondern}{summary}
\begin{itemize}
    \item \GermanText{\textbf{sondern} bedeutet: nicht A, sondern B.}
    \item \GermanText{Der erste Teil ist normalerweise verneint: \textbf{nicht}, \textbf{kein}, \textbf{nie}, \textbf{niemand}.}
    \item \GermanText{Vor \textbf{sondern} steht normalerweise ein Komma.}
    \item \GermanText{\textbf{sondern} verändert die Wortstellung nicht.}
    \item \GermanText{\textbf{aber} = Gegensatz; \textbf{sondern} = Korrektur nach Verneinung.}
\end{itemize}
\end{grammarentry}

\end{document}
